\documentclass{amsart}
% For dealing with references we use the comment environment
\usepackage{verbatim}
\newenvironment{reference}{\comment}{\endcomment}
%\newenvironment{reference}{}{}
\newenvironment{slogan}{\comment}{\endcomment}
\newenvironment{history}{\comment}{\endcomment}

% For commutative diagrams we use Xy-pic
\usepackage[all]{xy}

% We use 2cell for 2-commutative diagrams.
\xyoption{2cell}
\UseAllTwocells

% We use multicol for the list of chapters between chapters
\usepackage{multicol}

% This is generally recommended for better output
\usepackage{lmodern}
\usepackage[T1]{fontenc}

% For cross-file-references

% Package for hypertext links:
\usepackage{hyperref}

% For any local file, say "hello.tex" you want to link to please

% Theorem environments.
%
\theoremstyle{plain}
\newtheorem{theorem}[subsection]{Theorem}
\newtheorem{proposition}[subsection]{Proposition}
\newtheorem{lemma}[subsection]{Lemma}

\theoremstyle{definition}
\newtheorem{definition}[subsection]{Definition}
\newtheorem{example}[subsection]{Example}
\newtheorem{exercise}[subsection]{Exercise}
\newtheorem{situation}[subsection]{Situation}

\theoremstyle{remark}
\newtheorem{remark}[subsection]{Remark}
\newtheorem{remarks}[subsection]{Remarks}

\numberwithin{equation}{subsection}

% Macros
%
\def\lim{\mathop{\mathrm{lim}}\nolimits}
\def\colim{\mathop{\mathrm{colim}}\nolimits}
\def\Spec{\mathop{\mathrm{Spec}}}
\def\Hom{\mathop{\mathrm{Hom}}\nolimits}
\def\Ext{\mathop{\mathrm{Ext}}\nolimits}
\def\SheafHom{\mathop{\mathcal{H}\!\mathit{om}}\nolimits}
\def\SheafExt{\mathop{\mathcal{E}\!\mathit{xt}}\nolimits}
\def\Sch{\mathit{Sch}}
\def\Mor{\mathop{\mathrm{Mor}}\nolimits}
\def\Ob{\mathop{\mathrm{Ob}}\nolimits}
\def\Sh{\mathop{\mathit{Sh}}\nolimits}
\def\NL{\mathop{N\!L}\nolimits}
\def\CH{\mathop{\mathrm{CH}}\nolimits}
\def\proetale{{pro\text{-}\acute{e}tale}}
\def\etale{{\acute{e}tale}}
\def\QCoh{\mathit{QCoh}}
\def\Ker{\mathop{\mathrm{Ker}}}
\def\Im{\mathop{\mathrm{Im}}}
\def\Coker{\mathop{\mathrm{Coker}}}
\def\Coim{\mathop{\mathrm{Coim}}}

% Boxtimes
%
\DeclareMathSymbol{\boxtimes}{\mathbin}{AMSa}{"02}

%
% Macros for moduli stacks/spaces
%
\def\QCohstack{\mathcal{QC}\!\mathit{oh}}
\def\Cohstack{\mathcal{C}\!\mathit{oh}}
\def\Spacesstack{\mathcal{S}\!\mathit{paces}}
\def\Quotfunctor{\mathrm{Quot}}
\def\Hilbfunctor{\mathrm{Hilb}}
\def\Curvesstack{\mathcal{C}\!\mathit{urves}}
\def\Polarizedstack{\mathcal{P}\!\mathit{olarized}}
\def\Complexesstack{\mathcal{C}\!\mathit{omplexes}}
% \Pic is the operator that assigns to X its picard group, usage \Pic(X)
% \Picardstack_{X/B} denotes the Picard stack of X over B
% \Picardfunctor_{X/B} denotes the Picard functor of X over B
\def\Pic{\mathop{\mathrm{Pic}}\nolimits}
\def\Picardstack{\mathcal{P}\!\mathit{ic}}
\def\Picardfunctor{\mathrm{Pic}}
\def\Deformationcategory{\mathcal{D}\!\mathit{ef}}

\setlength{\emergencystretch}{3em}
\begin{document}
\title[Stacks Project, Tag 0BMA]{322 --- Purity of finite etale covers of regular local punctured spectra}
\maketitle
\noindent Copyright (C) 2005--2025 Johan de Jong. Stacks Project authors. Source: \href{https://stacks.math.columbia.edu/tag/0BMA}{Tag 0BMA}.

\noindent Editorial difficulty note (not author text). Difficulty H2: The dimension-two case requires proving flatness of a finite normal algebra through Cohen-Macaulay structure and then excluding a zero-dimensional branch locus by the discriminant. In higher dimension one must reduce covers across a regular parameter while retaining the extension property, completing a nontrivial dimension induction..

\noindent Editorial provenance note (not author text). History: excerpt from revision \texttt{a0444\allowbreak 6e57e\allowbreak c1fbc\allowbreak 252a8\allowbreak 71afc\allowbreak ec775\allowbreak 2fb28\allowbreak 07b14}, pione.tex, lines 5066--5115. Reference presentation was adapted; original-author context and core support blocks were retained or added. The mathematical wording is unchanged. Licensed under GNU FDL 1.2 or later; no invariant sections or cover texts. See ../licenses/COPYING. Context blocks reproduce author definitions and supporting results; they are not additional counted problems.

\begin{definition}
\label{morphisms-definition-unramified}
Let $f : X \to S$ be a morphism of schemes.
\begin{enumerate}
\item We say that $f$ is {\it unramified at $x \in X$} if
there exists an affine open neighbourhood $\Spec(A) = U \subset X$
of $x$ and affine open $\Spec(R) = V \subset S$
with $f(U) \subset V$ such that the induced ring map
$R \to A$ is unramified.
\item We say that $f$ is {\it G-unramified at $x \in X$} if
there exists an affine open neighbourhood $\Spec(A) = U \subset X$
of $x$ and affine open $\Spec(R) = V \subset S$
with $f(U) \subset V$ such that the induced ring map
$R \to A$ is G-unramified.
\item We say that $f$ is {\it unramified} if it is unramified
at every point of $X$.
\item We say that $f$ is {\it G-unramified} if it is G-unramified
at every point of $X$.
\end{enumerate}
\end{definition}

\begin{definition}
\label{morphisms-definition-etale}
Let $f : X \to S$ be a morphism of schemes.
\begin{enumerate}
\item We say that $f$ is {\it \'etale at $x \in X$} if
there exists an affine open neighbourhood $\Spec(A) = U \subset X$
of $x$ and affine open $\Spec(R) = V \subset S$
with $f(U) \subset V$ such that the induced ring map
$R \to A$ is \'etale.
\item We say that $f$ is {\it \'etale} if it is \'etale at every point of $X$.
\item A morphism of affine schemes $f : X \to S$ is called
{\it standard \'etale} if $X \to S$ is isomorphic to
$$
\Spec(R[x]_h/(g)) \to \Spec(R)
$$
where $R \to R[x]_h/(g)$ is a standard \'etale ring map, see
Algebra, Definition \ref{algebra-definition-standard-etale},
i.e., $g$ is monic and $g'$ invertible in $R[x]_h/(g)$.
\end{enumerate}
\end{definition}

\begin{situation}
\label{situation-local-lefschetz}
Let $(A, \mathfrak m)$ be a Noetherian local ring and $f \in \mathfrak m$.
We set $X = \Spec(A)$ and $X_0 = \Spec(A/fA)$ and we
let $U = X \setminus \{\mathfrak m\}$ and
$U_0 = X_0 \setminus \{\mathfrak m\}$ be the punctured spectrum of
$A$ and $A/fA$.
\end{situation}

\begin{definition}
\label{algebra-definition-standard-etale}
Let $R$ be a ring. Let $g , f  \in R[x]$.
Assume that $f$ is monic and the derivative $f'$ is invertible in
the localization $R[x]_g/(f)$.
In this case the ring map $R \to R[x]_g/(f)$ is said to be
{\it standard \'etale}.
\end{definition}

\begin{lemma}
\label{lemma-reformulate-purity-normal}
Let $(A, \mathfrak m)$ be a Noetherian local ring. Set $X = \Spec(A)$
and let $U = X \setminus \{\mathfrak m\}$. Assume $A$ is normal
of dimension $\geq 2$. The functor
$$
\textit{F\'Et}_U \longrightarrow
\left\{
\begin{matrix}
\text{finite normal }A\text{-algebras }B\text{ such} \\
\text{that }\Spec(B) \to X\text{ is \'etale over }U
\end{matrix}
\right\},
\quad
V \longmapsto \Gamma(V, \mathcal{O}_V)
$$
is an equivalence. Moreover, $V = Y \times_X U$ for some $Y \to X$
finite \'etale if and only if $B = \Gamma(V, \mathcal{O}_V)$
is finite \'etale over $A$.
\end{lemma}

\begin{proof}
Observe that $\text{depth}(A) \geq 2$ because $A$ is normal
(Serre's criterion for normality, Algebra, Lemma
\href{https://stacks.math.columbia.edu/tag/031S}{lemma criterion normal (Tag 031S)}).
Thus the final statement follows from Lemma \href{https://stacks.math.columbia.edu/tag/0BLK}{lemma reformulate purity (Tag 0BLK)}.
Given $\pi : V \to U$ finite \'etale, set $B = \Gamma(V, \mathcal{O}_V)$.
If we can show that $B$ is normal and finite over $A$, then
we obtain the displayed functor. Since there is an obvious
quasi-inverse functor, this is also all that we have to show.

\medskip\noindent
Since $A$ is normal, the scheme $V$ is normal
(Descent, Lemma \href{https://stacks.math.columbia.edu/tag/034F}{lemma normal local smooth (Tag 034F)}).
Hence $V$ is a finite disjoint union of integral schemes
(Properties, Lemma \href{https://stacks.math.columbia.edu/tag/033M}{lemma normal Noetherian (Tag 033M)}).
Thus we may assume $V$ is integral.
In this case the function field $L$ of $V$
(Morphisms, Section \href{https://stacks.math.columbia.edu/tag/01RR}{section rational maps (Tag 01RR)})
is a finite separable extension of the fraction field of $A$
(because we get it by looking at the generic fibre
of $V \to U$ and using Morphisms, Lemma
\href{https://stacks.math.columbia.edu/tag/02GL}{lemma etale over field (Tag 02GL)}).
By Algebra, Lemma
\href{https://stacks.math.columbia.edu/tag/032L}{lemma Noetherian normal domain finite separable extension (Tag 032L)}
the integral closure $B' \subset L$ of $A$ in $L$ is finite over $A$.
By More on Algebra, Lemma \href{https://stacks.math.columbia.edu/tag/0BM4}{lemma integral closure reflexive (Tag 0BM4)}
we see that $B'$ is a reflexive $A$-module, which in turn implies
that $\text{depth}_A(B') \geq 2$ by
More on Algebra, Lemma \href{https://stacks.math.columbia.edu/tag/0AVB}{lemma reflexive over normal (Tag 0AVB)}.

\medskip\noindent
Let $f \in \mathfrak m$. Then $B_f = \Gamma(V \times_U D(f), \mathcal{O}_V)$
(Properties, Lemma \href{https://stacks.math.columbia.edu/tag/01P7}{lemma invert f sections (Tag 01P7)}).
Hence $B'_f = B_f$ because $B_f$ is normal (see above),
finite over $A_f$ with fraction field $L$.
It follows that $V = \Spec(B') \times_X U$.
Then we conclude that $B = B'$ from
Lemma \ref{lemma-sections-over-punctured-spec}
applied to $\Spec(B') \to X$.
This lemma applies because the localizations $B'_{\mathfrak m'}$
of $B'$ at maximal ideals $\mathfrak m' \subset B'$ lying over
$\mathfrak m$ have depth $\geq 2$ by
Algebra, Lemma \href{https://stacks.math.columbia.edu/tag/0AUK}{lemma depth goes down finite (Tag 0AUK)}
and the remark on depth in the preceding paragraph.
\end{proof}


\begin{lemma}
\label{lemma-sections-over-punctured-spec}
Let $(A, \mathfrak m)$ be a Noetherian local ring. Set $X = \Spec(A)$
and let $U = X \setminus \{\mathfrak m\}$.
Let $\pi : Y \to X$ be a finite morphism such that
$\text{depth}(\mathcal{O}_{Y, y}) \geq 2$ for all closed points
$y \in Y$.
Then $Y$ is the spectrum of $B = \mathcal{O}_Y(\pi^{-1}(U))$.
\end{lemma}

\begin{proof}
Set $V = \pi^{-1}(U)$ and denote $\pi' : V \to U$ the restriction of $\pi$.
Consider the $\mathcal{O}_X$-module map
$$
\pi_*\mathcal{O}_Y \longrightarrow j_*\pi'_*\mathcal{O}_V
$$
where $j : U \to X$ is the inclusion morphism. We claim
Divisors, Lemma \href{https://stacks.math.columbia.edu/tag/0E9I}{lemma depth 2 hartog (Tag 0E9I)}
applies to this map. If so, then $B = \Gamma(Y, \mathcal{O}_Y)$
and we see that the lemma holds. Let $x \in X$ be the closed point.
It suffices to show that
$\text{depth}((\pi_*\mathcal{O}_Y)_x) \geq 2$.
Let $y_1, \ldots, y_n \in Y$ be the points mapping to $x$.
By Algebra, Lemma \href{https://stacks.math.columbia.edu/tag/0AUK}{lemma depth goes down finite (Tag 0AUK)}
it suffices to show that
$\text{depth}(\mathcal{O}_{Y, y_i}) \geq 2$ for $i = 1, \ldots, n$.
Since this is the assumption of the lemma the proof is complete.
\end{proof}


\begin{lemma}
\label{lemma-lift-purity}
In Situation \ref{situation-local-lefschetz}.
Let $V$ be finite \'etale over $U$. Assume
\begin{enumerate}
\item $A$ has depth $\geq 3$,
\item $V_0 = V \times_U U_0$ is equal to $Y_0 \times_{X_0} U_0$
for some $Y_0 \to X_0$ finite \'etale.
\end{enumerate}
Then $V = Y \times_X U$ for some $Y \to X$ finite \'etale.
\end{lemma}

\begin{proof}
The assumption of depth forces
$H^1_\mathfrak m(A) = H^2_\mathfrak m(A) = 0$, see
Dualizing Complexes, Lemma \href{https://stacks.math.columbia.edu/tag/0AVZ}{lemma depth (Tag 0AVZ)}.
Hence Lemma \ref{lemma-lift-purity-general} applies.
\end{proof}


\begin{lemma}
\label{lemma-lift-purity-general}
In Situation \ref{situation-local-lefschetz}.
Let $V$ be finite \'etale over $U$. Assume
\begin{enumerate}
\item $H^1_\mathfrak m(A)$ and $H^2_\mathfrak m(A)$
are annihilated by a power of $f$,
\item $V_0 = V \times_U U_0$ is equal to $Y_0 \times_{X_0} U_0$
for some $Y_0 \to X_0$ finite \'etale.
\end{enumerate}
Then $V = Y \times_X U$ for some $Y \to X$ finite \'etale.
\end{lemma}

\begin{proof}
We reduce to the complete case using Lemma \ref{lemma-purity-and-completion}.
(The assumptions carry over; use Dualizing Complexes, Lemma
\href{https://stacks.math.columbia.edu/tag/0ALZ}{lemma torsion change rings (Tag 0ALZ)}.)

\medskip\noindent
In the complete case we can lift $Y_0 \to X_0$ to a finite \'etale
morphism $Y \to X$ by
More on Algebra, Lemma \ref{more-algebra-lemma-finite-etale-equivalence};
observe that $(A, fA)$ is a henselian pair by
More on Algebra, Lemma \href{https://stacks.math.columbia.edu/tag/0ALJ}{lemma complete henselian (Tag 0ALJ)}.
Then we can use Lemma \ref{lemma-fully-faithful-minimal}
to see that $V$ is isomorphic to $Y \times_X U$ and
the proof is complete.
\end{proof}


\begin{lemma}
\label{lemma-purity-and-completion}
Let $(A, \mathfrak m)$ be a Noetherian local ring. Set $X = \Spec(A)$
and let $U = X \setminus \{\mathfrak m\}$.
Let $V$ be finite \'etale over $U$.
Let $A^\wedge$ be the $\mathfrak m$-adic completion of $A$,
let $X' = \Spec(A^\wedge)$ and let $U'$ and $V'$ be the base changes of
$U$ and $V$ to $X'$. The following are equivalent
\begin{enumerate}
\item $V = Y \times_X U$ for some $Y \to X$ finite \'etale, and
\item $V' = Y' \times_{X'} U'$ for some $Y' \to X'$ finite \'etale.
\end{enumerate}
\end{lemma}

\begin{proof}
The implication (1) $\Rightarrow$ (2) follows from taking the base change
of a solution $Y \to X$. Let $Y' \to X'$ be as in (2).
By Lemma \ref{lemma-fill-in-missing} we can find $Y \to X$ finite
such that $V = Y \times_X U$ and $Y' = Y \times_X X'$.
By descent we see that $Y \to X$ is finite \'etale
(Algebra, Lemmas \href{https://stacks.math.columbia.edu/tag/03C4}{lemma descend properties modules (Tag 03C4)} and
\href{https://stacks.math.columbia.edu/tag/00U2}{lemma etale (Tag 00U2)}). This finishes the proof.
\end{proof}


\begin{lemma}
\label{lemma-fill-in-missing}
In Situation \ref{situation-local-lefschetz}. Let $V \to U$ be a finite
morphism.  Let $A^\wedge$ be the $\mathfrak m$-adic completion of $A$,
let $X' = \Spec(A^\wedge)$ and let $U'$ and $V'$ be the base changes of
$U$ and $V$ to $X'$. If $Y' \to X'$ is a finite morphism such that
$V' = Y' \times_{X'} U'$, then there exists a finite morphism $Y \to X$
such that $V = Y \times_X U$ and $Y' = Y \times_X X'$.
\end{lemma}

\begin{proof}
This is a straightforward application of
More on Algebra, Proposition \href{https://stacks.math.columbia.edu/tag/05ER}{proposition equivalence (Tag 05ER)}.
Namely, choose generators $f_1, \ldots, f_t$ of $\mathfrak m$.
For each $i$ write $V \times_U D(f_i) = \Spec(B_i)$.
For $1 \leq i, j \leq n$ we obtain an isomorphism
$\alpha_{ij} : (B_i)_{f_j} \to (B_j)_{f_i}$ of $A_{f_if_j}$-algebras
because the spectrum of both represent $V \times_U D(f_if_j)$.
Write $Y' = \Spec(B')$. Since $V \times_U U' = Y \times_{X'} U'$
we get isomorphisms $\alpha_i : B'_{f_i} \to B_i \otimes_A A^\wedge$.
A straightforward argument shows that $(B', B_i, \alpha_i, \alpha_{ij})$
is an object of $\text{Glue}(A \to A^\wedge, f_1, \ldots, f_t)$, see
More on Algebra, Remark \href{https://stacks.math.columbia.edu/tag/05EL}{remark glueing data (Tag 05EL)}.
Applying the proposition cited above (and using
More on Algebra, Remark \href{https://stacks.math.columbia.edu/tag/05EU}{remark formal glueing algebras (Tag 05EU)}
to obtain the algebra structure) we find an $A$-algebra $B$ such that
$\text{Can}(B)$ is isomorphic to $(B', B_i, \alpha_i, \alpha_{ij})$.
Setting $Y = \Spec(B)$ we see that $Y \to X$ is a morphism
which comes equipped with compatible isomorphisms
$V \cong Y \times_X U$ and $Y' = Y \times_X X'$ as desired.
\end{proof}


\begin{lemma}
\label{lemma-fully-faithful-minimal}
\begin{reference}
\href{https://stacks.math.columbia.edu/bibliography}{Bibliography: Bhatt-local (Corollary 1.11)}
\end{reference}
In Situation \ref{situation-local-lefschetz}. Assume
\begin{enumerate}
\item $H^1_\mathfrak m(A)$ and $H^2_\mathfrak m(A)$ are
annihilated by a power of $f$, and
\item $A$ is henselian or more generally $(A, (f))$ is a henselian pair.
\end{enumerate}
Then the restriction functor
$\textit{F\'Et}_U \longrightarrow \textit{F\'Et}_{U_0}$
is fully faithful.
\end{lemma}

\begin{proof}
By Lemma \href{https://stacks.math.columbia.edu/tag/0BLI}{lemma fully faithful henselian completion (Tag 0BLI)}
we may assume that $A$ is a Noetherian complete local ring.
(The assumptions carry over; use
Dualizing Complexes, Lemma \href{https://stacks.math.columbia.edu/tag/0ALZ}{lemma torsion change rings (Tag 0ALZ)}.)
By Lemma \ref{lemma-restriction-fully-faithful}
the result follows from
Algebraic and Formal Geometry, Lemma
\ref{algebraization-lemma-fully-faithful-alternative}.
\end{proof}


\begin{lemma}
\label{lemma-restriction-fully-faithful}
Let $X$ be a Noetherian scheme and let $Y \subset X$ be a closed subscheme
with ideal sheaf $\mathcal{I} \subset \mathcal{O}_X$.
Assume the completion functor
$$
\textit{Coh}(\mathcal{O}_X)
\longrightarrow 
\textit{Coh}(X, \mathcal{I}),\quad
\mathcal{F} \longmapsto \mathcal{F}^\wedge
$$
is fully faithful on the full subcategory of finite locally free objects
(see above).
Then the restriction functor $\textit{F\'Et}_X \to \textit{F\'Et}_Y$
is fully faithful.
\end{lemma}

\begin{proof}
Since the category of finite \'etale coverings has an
internal hom (Lemma \ref{lemma-internal-hom-finite-etale})
it suffices to prove the following: Given $U$ finite \'etale over $X$
and a morphism $t : Y \to U$ over $X$ there exists a unique section
$s : X \to U$ such that $t = s|_Y$. Picture
$$
\xymatrix{
& U \ar[d]^f \\
Y \ar[r] \ar[ru] & X \ar@{..>}@/^1em/[u]
}
$$
Finding the dotted arrow $s$ is the same thing as finding an
$\mathcal{O}_X$-algebra map
$$
s^\sharp : f_*\mathcal{O}_U \longrightarrow \mathcal{O}_X
$$
which reduces modulo the ideal sheaf of $Y$ to the given algebra map
$t^\sharp : f_*\mathcal{O}_U \to \mathcal{O}_Y$.
By Lemma \ref{lemma-thickening} we can lift $t$ uniquely to a compatible
system of maps $t_n : Y_n \to U$ and hence a map
$$
\lim t_n^\sharp : f_*\mathcal{O}_U \longrightarrow \lim \mathcal{O}_{Y_n}
$$
of sheaves of algebras on $X$.
Since $f_*\mathcal{O}_U$ is a finite locally free $\mathcal{O}_X$-module,
we conclude that we get a unique $\mathcal{O}_X$-module map
$\sigma : f_*\mathcal{O}_U \to \mathcal{O}_X$ whose completion
is $\lim t_n^\sharp$. To see that $\sigma$ is an algebra homomorphism,
we need to check that the diagram
$$
\xymatrix{
f_*\mathcal{O}_U \otimes_{\mathcal{O}_X} f_*\mathcal{O}_U
\ar[r] \ar[d]_{\sigma \otimes \sigma} &
f_*\mathcal{O}_U \ar[d]^\sigma \\
\mathcal{O}_X \otimes_{\mathcal{O}_X} \mathcal{O}_X \ar[r] &
\mathcal{O}_X
}
$$
commutes. For every $n$ we know this diagram commutes after restricting
to $Y_n$, i.e., the diagram commutes after applying the completion functor.
Hence by faithfulness of the completion functor we conclude.
\end{proof}


\begin{lemma}
\label{algebraization-lemma-fully-faithful-alternative}
Let $A$ be a Noetherian ring. Let $f \in \mathfrak a \subset A$
be an element of an ideal of $A$. Let $U = \Spec(A) \setminus V(\mathfrak a)$.
Assume
\begin{enumerate}
\item $A$ is $f$-adically complete,
\item $H^1_\mathfrak a(A)$ and $H^2_\mathfrak a(A)$ are
annihilated by a power of $f$.
\end{enumerate}
Then the completion functor
$$
\textit{Coh}(\mathcal{O}_U)
\longrightarrow
\textit{Coh}(U, I\mathcal{O}_U),
\quad
\mathcal{F} \longmapsto \mathcal{F}^\wedge
$$
is fully faithful on the full subcategory of
finite locally free objects.
\end{lemma}

\begin{proof}
By Lemma \ref{algebraization-lemma-completion-fully-faithful}
it suffices to show that
$$
\Gamma(U, \mathcal{O}_U) =
\lim \Gamma(U, \mathcal{O}_U/I^n\mathcal{O}_U)
$$
This follows immediately from
Lemma \ref{algebraization-lemma-alternative-H0}.
\end{proof}


\begin{lemma}
\label{algebraization-lemma-alternative-H0}
Let $A$ be a Noetherian ring. Let $f \in \mathfrak a \subset A$
be an element of an ideal of $A$. Let $M$ be a finite $A$-module.
Assume
\begin{enumerate}
\item $A$ is $f$-adically complete,
\item $H^1_\mathfrak a(M)$ and $H^2_\mathfrak a(M)$ are
annihilated by a power of $f$.
\end{enumerate}
Then with $U = \Spec(A) \setminus V(\mathfrak a)$ the map
$$
\Gamma(U, \widetilde{M})
\longrightarrow
\lim \Gamma(U, \widetilde{M/f^nM})
$$
is an isomorphism.
\end{lemma}

\begin{proof}
We may apply Lemma \href{https://stacks.math.columbia.edu/tag/0BLD}{lemma formal functions principal (Tag 0BLD)} to $U$ and
$\mathcal{F} = \widetilde{M}|_U$ because $\mathcal{F}$ is a Noetherian object in
the category of coherent $\mathcal{O}_U$-modules.
Since $H^1(U, \mathcal{F}) = H^2_\mathfrak a(M)$
(Local Cohomology, Lemma
\href{https://stacks.math.columbia.edu/tag/0BK0}{lemma finiteness pushforwards and H1 local (Tag 0BK0)})
is annihilated by a power of $f$, we see that
its $f$-adic Tate module is zero.
Hence the lemma shows $\lim H^0(U, \mathcal{F}/f^n \mathcal{F})$
is equal to the usual $f$-adic completion of $H^0(U, \mathcal{F})$.
Consider the short exact sequence
$$
0 \to M/H^0_\mathfrak a(M) \to H^0(U, \mathcal{F}) \to H^1_\mathfrak a(M) \to 0
$$
of Local Cohomology, Lemma
\href{https://stacks.math.columbia.edu/tag/0BK0}{lemma finiteness pushforwards and H1 local (Tag 0BK0)}.
Since $M/H^0_\mathfrak a(M)$ is a finite $A$-module, it is
complete, see Algebra, Lemma \href{https://stacks.math.columbia.edu/tag/00MA}{lemma completion tensor (Tag 00MA)}.
Since $H^1_\mathfrak a(M)$ is killed by a power of $f$, we conclude
from Algebra, Lemma \href{https://stacks.math.columbia.edu/tag/0BNG}{lemma completion differ by torsion (Tag 0BNG)}
that $H^0(U, \mathcal{F})$ is complete as well. This finishes the proof.
\end{proof}


\begin{lemma}
\label{algebraization-lemma-completion-fully-faithful}
Let $X$ be a Noetherian scheme and let $Y \subset X$ be a closed subscheme.
Let $Y_n \subset X$ be the $n$th infinitesimal neighbourhood of $Y$ in $X$.
Consider the following conditions
\begin{enumerate}
\item $X$ is quasi-affine and
$\Gamma(X, \mathcal{O}_X) \to \lim \Gamma(Y_n, \mathcal{O}_{Y_n})$
is an isomorphism,
\item $X$ has an ample invertible module $\mathcal{L}$ and
$\Gamma(X, \mathcal{L}^{\otimes m}) \to
\lim \Gamma(Y_n, \mathcal{L}^{\otimes m}|_{Y_n})$
is an isomorphism for all $m \gg 0$,
\item for every finite locally free $\mathcal{O}_X$-module
$\mathcal{E}$ the map
$\Gamma(X, \mathcal{E}) \to \lim \Gamma(Y_n, \mathcal{E}|_{Y_n})$
is an isomorphism, and
\item the completion functor
$\textit{Coh}(\mathcal{O}_X) \to \textit{Coh}(X, \mathcal{I})$
is fully faithful on the full subcategory of finite locally free
objects.
\end{enumerate}
Then (1) $\Rightarrow$ (2) $\Rightarrow$ (3) $\Rightarrow$ (4)
and (4) $\Rightarrow$ (3).
\end{lemma}

\begin{proof}
Proof of (3) $\Rightarrow$ (4). If $\mathcal{F}$ and $\mathcal{G}$
are finite locally free on $X$, then considering
$\mathcal{H} = \SheafHom_{\mathcal{O}_X}(\mathcal{G}, \mathcal{F})$
and using Cohomology of Schemes, Lemma
\href{https://stacks.math.columbia.edu/tag/0882}{lemma completion internal hom (Tag 0882)}
we see that (3) implies (4).

\medskip\noindent
Proof of (2) $\rightarrow$ (3). Namely, let $\mathcal{L}$ be ample
on $X$ and suppose that $\mathcal{E}$ is a
finite locally free $\mathcal{O}_X$-module.
We claim we can find a universally exact sequence
$$
0 \to \mathcal{E} \to
(\mathcal{L}^{\otimes p})^{\oplus r} \to
(\mathcal{L}^{\otimes q})^{\oplus s}
$$
for some $r, s \geq 0$ and $0 \ll p \ll q$. If this holds, then
using the exact sequence
$$
0 \to \lim \Gamma(\mathcal{E}|_{Y_n}) \to
\lim \Gamma((\mathcal{L}^{\otimes p})^{\oplus r}|_{Y_n}) \to
\lim \Gamma((\mathcal{L}^{\otimes q})^{\oplus s}|_{Y_n})
$$
and the isomorphisms in (2) we get the isomorphism in (3).
To prove the claim, consider the dual locally free module
$\SheafHom_{\mathcal{O}_X}(\mathcal{E}, \mathcal{O}_X)$
and apply
Properties, Proposition \href{https://stacks.math.columbia.edu/tag/01Q3}{proposition characterize ample (Tag 01Q3)}
to find a surjection
$$
(\mathcal{L}^{\otimes -p})^{\oplus r}
\longrightarrow
\SheafHom_{\mathcal{O}_X}(\mathcal{E}, \mathcal{O}_X)
$$
Taking duals we obtain the first map in the exact sequence
(it is universally injective because being a surjection is universal).
Repeat with the cokernel to get the second. Some details omitted.

\medskip\noindent
Proof of (1) $\Rightarrow$ (2). This is true because if $X$ is quasi-affine
then $\mathcal{O}_X$ is an ample invertible module, see
Properties, Lemma \href{https://stacks.math.columbia.edu/tag/01QE}{lemma quasi affine O ample (Tag 01QE)}.

\medskip\noindent
We omit the proof of (4) $\Rightarrow$ (3).
\end{proof}


\begin{lemma}
\label{lemma-ramification-quasi-finite-flat}
Let $f : X \to Y$ be a morphism of locally Noetherian schemes.
Let $x \in X$. Assume
\begin{enumerate}
\item $f$ is flat,
\item $f$ is quasi-finite at $x$,
\item $x$ is not a generic point of an irreducible component of $X$,
\item for specializations $x' \leadsto x$ with
$\dim(\mathcal{O}_{X, x'}) = 1$ our $f$ is unramified at $x'$.
\end{enumerate}
Then $f$ is \'etale at $x$.
\end{lemma}

\begin{proof}
Observe that the set of points where $f$ is unramified is the same as
the set of points where $f$ is \'etale and that this set is open.
See Morphisms, Definitions \ref{morphisms-definition-unramified}
and \ref{morphisms-definition-etale} and
Lemma \href{https://stacks.math.columbia.edu/tag/02GV}{lemma flat unramified etale (Tag 02GV)}.
To check $f$ is \'etale at $x$ we may work \'etale
locally on the base and on the
target (Descent, Lemmas \href{https://stacks.math.columbia.edu/tag/02VN}{lemma descending property etale (Tag 02VN)} and
\href{https://stacks.math.columbia.edu/tag/036W}{lemma etale etale local source (Tag 036W)}).
Thus we can apply More on Morphisms, Lemma
\href{https://stacks.math.columbia.edu/tag/02LK}{lemma etale makes quasi finite finite at point (Tag 02LK)}
and assume that $f : X \to Y$ is finite and that $x$ is the unique
point of $X$ lying over $y = f(x)$.
Then it follows that $f$ is finite locally free
(Morphisms, Lemma \href{https://stacks.math.columbia.edu/tag/02KB}{lemma finite flat (Tag 02KB)}).

\medskip\noindent
Assume $f$ is finite locally free and that $x$ is the unique point of
$X$ lying over $y = f(x)$. By
Discriminants, Lemma \href{https://stacks.math.columbia.edu/tag/0BJF}{lemma discriminant (Tag 0BJF)}
we find a locally principal closed subscheme $D_\pi \subset Y$
such that $y' \in D_\pi$ if and only if there exists an $x' \in X$
with $f(x') = y'$ and $f$ ramified at $x'$. Thus we have to prove
that $y \not \in D_\pi$. Assume $y \in D_\pi$ to get a contradiction.

\medskip\noindent
By condition (3) we have $\dim(\mathcal{O}_{X, x}) \geq 1$.
We have $\dim(\mathcal{O}_{X, x}) = \dim(\mathcal{O}_{Y, y})$ by
Algebra, Lemma \href{https://stacks.math.columbia.edu/tag/00ON}{lemma dimension base fibre equals total (Tag 00ON)}.
By Lemma \href{https://stacks.math.columbia.edu/tag/0BJG}{lemma find point codim 1 (Tag 0BJG)}
we can find $y' \in D_\pi$ specializing to $y$
with $\dim(\mathcal{O}_{Y, y'}) = 1$.
Choose $x' \in X$ with $f(x') = y'$ where $f$ is ramified. Since $f$
is finite it is closed, and hence $x' \leadsto x$.
We have $\dim(\mathcal{O}_{X, x'}) = \dim(\mathcal{O}_{Y, y'}) = 1$
as before. This contradicts property (4).
\end{proof}


\begin{lemma}
\label{lemma-internal-hom-finite-etale}
Let $X$ be a scheme. Given $U, V$ finite \'etale over $X$ there
exists a scheme $W$ finite \'etale over $X$ such that
$$
\Mor_X(X, W) = \Mor_X(U, V)
$$
and such that the same remains true after any base change.
\end{lemma}

\begin{proof}
By More on Morphisms, Lemma
\href{https://stacks.math.columbia.edu/tag/0BL5}{lemma hom from finite locally free separated lqf (Tag 0BL5)}
there exists a scheme $W$ representing $\mathit{Mor}_X(U, V)$.
(Use that an \'etale morphism is locally quasi-finite by
Morphisms, Lemmas \href{https://stacks.math.columbia.edu/tag/03WS}{lemma etale locally quasi finite (Tag 03WS)}
and that a finite morphism is separated.)
This scheme clearly satisfies the formula after any base change.
To finish the proof we have to show that $W \to X$ is finite \'etale.
This we may do after replacing $X$ by the members of an \'etale
covering (Descent, Lemmas \href{https://stacks.math.columbia.edu/tag/02LA}{lemma descending property finite (Tag 02LA)}
and \href{https://stacks.math.columbia.edu/tag/02KU}{lemma descending property separated (Tag 02KU)}).
Thus by \'Etale Morphisms, Lemma \href{https://stacks.math.columbia.edu/tag/04HN}{lemma finite etale etale local (Tag 04HN)}
we may assume that $U = \coprod_{i = 1, \ldots, n} X$
and $V = \coprod_{j = 1, \ldots, m} X$.
In this case
$W = \coprod_{\alpha : \{1, \ldots, n\} \to \{1, \ldots, m\}} X$
by inspection (details omitted) and the proof is complete.
\end{proof}


\begin{lemma}
\label{lemma-thickening}
Let $X \subset X'$ be a thickening of schemes. The functor
$$
\textit{F\'Et}_{X'} \longrightarrow \textit{F\'Et}_X,\quad
U' \longmapsto U' \times_{X'} X
$$
is an equivalence of categories.
\end{lemma}

\begin{proof}
For a discussion of thickenings see
More on Morphisms, Section \href{https://stacks.math.columbia.edu/tag/04EW}{section thickenings (Tag 04EW)}.
Let $U' \to X'$ be an \'etale morphism such that $U = U' \times_{X'} X \to X$
is finite \'etale. Then $U' \to X'$ is finite \'etale as well.
This follows for example from More on Morphisms, Lemma
\href{https://stacks.math.columbia.edu/tag/0BPG}{lemma properties that extend over thickenings (Tag 0BPG)}.
Now, if $X \subset X'$ is a finite order thickening then this remark
combined with \'Etale Morphisms, Theorem
\href{https://stacks.math.columbia.edu/tag/039R}{theorem remarkable equivalence (Tag 039R)}
proves the lemma. Below we will prove the lemma for general thickenings, but
we suggest the reader skip the proof.

\medskip\noindent
Let $X' = \bigcup X_i'$ be an affine open covering. Set
$X_i = X \times_{X'} X_i'$, $X_{ij}' = X'_i \cap X'_j$,
$X_{ij} = X \times_{X'} X_{ij}'$, $X_{ijk}' = X'_i \cap X'_j \cap X'_k$,
$X_{ijk} = X \times_{X'} X_{ijk}'$.
Suppose that we can prove
the theorem for each of the thickenings
$X_i \subset X'_i$, $X_{ij} \subset X_{ij}'$, and $X_{ijk} \subset X_{ijk}'$.
Then the result follows for $X \subset X'$ by relative glueing of
schemes, see
Constructions, Section \href{https://stacks.math.columbia.edu/tag/01LG}{section relative glueing (Tag 01LG)}.
Observe that the schemes $X_i'$, $X_{ij}'$, $X_{ijk}'$ are
each separated as open subschemes of affine schemes. Repeating the
argument one more time we reduce to the case where the schemes
$X'_i$, $X_{ij}'$, $X_{ijk}'$ are affine.

\medskip\noindent
In the affine case we have $X' = \Spec(A')$ and $X = \Spec(A'/I')$
where $I'$ is a locally nilpotent ideal. Then $(A', I')$ is a
henselian pair (More on Algebra, Lemma
\href{https://stacks.math.columbia.edu/tag/0ALI}{lemma locally nilpotent henselian (Tag 0ALI)})
and the result follows from Lemma \href{https://stacks.math.columbia.edu/tag/09ZS}{lemma gabber (Tag 09ZS)} (which is
much easier in this case).
\end{proof}


\begin{lemma}
\label{more-algebra-lemma-finite-etale-equivalence}
Let $(A, I)$ be a henselian pair. The functor $B \to B/IB$ determines
an equivalence between finite \'etale $A$-algebras and finite \'etale
$A/I$-algebras.
\end{lemma}

\begin{proof}
Let $B, B'$ be two $A$-algebras finite \'etale over $A$.
Then $B' \to B'' = B \otimes_A B'$ is finite \'etale as well
(Algebra, Lemmas \href{https://stacks.math.columbia.edu/tag/00U2}{lemma etale (Tag 00U2)} and
\href{https://stacks.math.columbia.edu/tag/02JK}{lemma base change integral (Tag 02JK)}).
Now we have $1$-to-$1$ correspondences between
\begin{enumerate}
\item $A$-algebra maps $B \to B'$,
\item sections of $B' \to B''$, and
\item idempotents $e$ of $B''$ such that $B' \to B'' \to eB''$ is
an isomorphism.
\end{enumerate}
The bijection between (2) and (3) sends $\sigma : B'' \to B'$
to $e$ such that $(1 - e)$ is the idempotent
that generates the kernel of $\sigma$ which exists by
Algebra, Lemmas \href{https://stacks.math.columbia.edu/tag/00U7}{lemma map between etale (Tag 00U7)} and
\href{https://stacks.math.columbia.edu/tag/00U8}{lemma surjective flat finitely presented (Tag 00U8)}.
There is a similar correspondence between
$A/I$-algebra maps $B/IB \to B'/IB'$ and idempotents
$\overline{e}$ of $B''/IB''$ such that
$B'/IB' \to B''/IB'' \to \overline{e}(B''/IB'')$ is
an isomorphism. However every idempotent $\overline{e}$ of $B''/IB''$
lifts uniquely to an idempotent $e$ of $B''$
(Lemma \href{https://stacks.math.columbia.edu/tag/09XI}{lemma characterize henselian pair (Tag 09XI)}).
Moreover, if $B'/IB' \to \overline{e}(B''/IB'')$ is an isomorphism,
then $B' \to eB''$ is an isomorphism too by Nakayama's lemma
(Algebra, Lemma \href{https://stacks.math.columbia.edu/tag/00DV}{lemma NAK (Tag 00DV)}).
In this way we see that the functor is fully faithful.

\medskip\noindent
Essential surjectivity. Let $A/I \to C$ be a finite \'etale map.
By Algebra, Lemma \href{https://stacks.math.columbia.edu/tag/04D1}{lemma lift etale (Tag 04D1)}
there exists an \'etale map $A \to B$ such that $B/IB \cong C$.
Let $B'$ be the integral closure of $A$ in $B$. 
By Lemma \href{https://stacks.math.columbia.edu/tag/09XH}{lemma helper finite type (Tag 09XH)} we have
$B'/IB' = C \times C'$ for some ring $C'$
and $B'_g \cong B_g$ for some $g \in B'$ mapping to $(1, 0) \in C \times C'$.
Since idempotents lift
(Lemma \href{https://stacks.math.columbia.edu/tag/09XI}{lemma characterize henselian pair (Tag 09XI)})
we get $B' = B'_1 \times B'_2$ with $C = B'_1/IB'_1$ and $C' = B'_2/IB'_2$.
The image of $g$ in $B'_1$ is invertible.
Then $B_g = B'_g = B'_1 \times (B_2)_g$ and this implies
that $A \to B'_1$ is \'etale.
We conclude that $B'_1$ is finite \'etale over $A$
(integral \'etale implies finite \'etale by
Algebra, Lemma \href{https://stacks.math.columbia.edu/tag/02JJ}{lemma characterize finite in terms of integral (Tag 02JJ)}
for example)
and the proof is done.
\end{proof}


\begin{lemma}
\label{lemma-local-purity}
Let $(A, \mathfrak m)$ be a regular local ring of dimension $d \geq 2$.
Set $X = \Spec(A)$ and $U = X \setminus \{\mathfrak m\}$. Then
\begin{enumerate}
\item the functor $\textit{F\'Et}_X \to \textit{F\'Et}_U$
is essentially surjective, i.e., purity holds for $A$,
\item any finite $A \to B$ with $B$ normal which
induces a finite \'etale morphism on punctured spectra is \'etale.
\end{enumerate}
\end{lemma}

\begin{proof}
Recall that a regular local ring is normal by
Algebra, Lemma \href{https://stacks.math.columbia.edu/tag/0567}{lemma regular normal (Tag 0567)}.
Hence (1) and (2) are equivalent by
Lemma \ref{lemma-reformulate-purity-normal}.
We prove the lemma by induction on $d$.

\medskip\noindent
The case $d = 2$. In this case $A \to B$ is flat.
Namely, we have going down for $A \to B$ by
Algebra, Proposition \href{https://stacks.math.columbia.edu/tag/00H8}{proposition going down normal integral (Tag 00H8)}.
Then $\dim(B_{\mathfrak m'}) = 2$ for all maximal ideals
$\mathfrak m' \subset B$ by
Algebra, Lemma \href{https://stacks.math.columbia.edu/tag/00ON}{lemma dimension base fibre equals total (Tag 00ON)}.
Then $B_{\mathfrak m'}$ is Cohen-Macaulay by
Algebra, Lemma \href{https://stacks.math.columbia.edu/tag/031S}{lemma criterion normal (Tag 031S)}.
Hence and this is the important step
Algebra, Lemma \href{https://stacks.math.columbia.edu/tag/00R4}{lemma CM over regular flat (Tag 00R4)}
applies to show $A \to B_{\mathfrak m'}$ is flat.
Then Algebra, Lemma \href{https://stacks.math.columbia.edu/tag/00HT}{lemma flat localization (Tag 00HT)}
shows $A \to B$ is flat. Thus we can apply
Lemma \ref{lemma-ramification-quasi-finite-flat}
(or you can directly argue using the easier
Discriminants, Lemma \href{https://stacks.math.columbia.edu/tag/0BJF}{lemma discriminant (Tag 0BJF)})
to see that $A \to B$ is \'etale.

\medskip\noindent
The case $d \geq 3$. Let $V \to U$ be finite \'etale.
Let $f \in \mathfrak m_A$, $f \not \in \mathfrak m_A^2$.
Then $A/fA$ is a regular local ring of dimension $d - 1 \geq 2$, see
Algebra, Lemma \href{https://stacks.math.columbia.edu/tag/00NQ}{lemma regular ring CM (Tag 00NQ)}.
Let $U_0$ be the punctured spectrum of $A/fA$ and let
$V_0 = V \times_U U_0$.
By Lemma \ref{lemma-lift-purity}
it suffices to show that $V_0$ is in the essential
image of $\textit{F\'Et}_{\Spec(A/fA)} \to \textit{F\'Et}_{U_0}$.
This follows from the induction hypothesis.
\end{proof}
\end{document}
